% ==========================================================================================
% PREAMBLE
% ==========================================================================================
%----------------------------------------------------------------------------------------
% GENERAL  - PACKAGES
%----------------------------------------------------------------------------------------
% GENERAL
\usepackage{etex} %Because of many packages --> Extended TeX.
\usepackage[utf8]{inputenc} %Due to vowels.
% \usepackage[british]{babel} %Define the language style.

%Load some mathematical packages.
\usepackage{amsmath}
\usepackage{amsfonts}
\usepackage{amsmath}
\usepackage{amssymb}
\usepackage{mathtools}
\usepackage{breqn} % breaking of equations (be careful with using ENDFLOAT  with this)

%  LAYOUT/PAGE/TEXT FORMATTING
\usepackage[left=1in, right=1in]{geometry} %Helps to structure the paper layout.
\usepackage{setspace} %Use double spacing.

\usepackage[Lenny]{styles/fncychap} %Design of the thesis. = Fancy chapter

\usepackage{fancyhdr} %To customize the headers and footers.
\usepackage[hang,bottom,stable,multiple]{footmisc} %Style of footnotes.
\usepackage{dsfont} %Nice style for the indicator function.

\usepackage[svgnames]{xcolor} % Enabling mixing colors and color's call by 'svgnames'
% define new colors (not limited to text obviously)
\definecolor{MyColor1}{rgb}{0.2,0.4,0.6} % mix personal color


% FLOATS
\usepackage{booktabs} %In case you need \cmidrule or \addlinespace in tables.
\usepackage{array} %To create tables and matrices.
\usepackage{hhline}
\usepackage{rotating} %To rotate a table/figure. e.g. \sidewaystable
\usepackage{tabularx} %An extended version of tabular.
\usepackage{float} %Allows for the 'H' option

\usepackage{graphicx} %For the graphics


\usepackage[margin=10pt, font=small, labelfont=bf, labelsep=endash]{caption} %Customize the captions.


% OTHER ENVIRONMENTS


\usepackage{textcomp}

\usepackage{amsthm} %For theorems, definitions etc.
\usepackage{thmtools} %For theorems, definitions etc.

\usepackage{appendix} %For the \appendixpage command.
\usepackage{etoolbox} %To remove the page number on \appendixpage.
\makeatletter %Remove page number on \appendixpage.
\patchcmd{\@chap@pppage}{\thispagestyle{plain}}{\thispagestyle{empty}}{}{}
\makeatother


\usepackage{styles/mcode} %To implement Matlab code.
\usepackage{listings} % For including code in your pdf

% VARIA  - does not mean useless
\usepackage{epstopdf} %For inserting .eps files into the document.
\usepackage{lipsum} %For the \lipsum command to generate a text.
\usepackage{datetime} %For the specification of the date.
\usepackage{chngcntr} %To use counterwithout.
\usepackage{xparse} %Load for \NewDocumentCommand command.
\usepackage{arydshln} %Due to the capability to draw horizontal/vertical dash-lines.

%  REFERENCING
% BIBLIOGRAPHY - BIBTEX
% \usepackage[sort,round]{natbib} %For the bibliography.
% \bibliographystyle{abbrvnat} %Reference style.

% BIBLIOGRAPHY - BIBLATEX
\usepackage[
backend=biber,
style=apa,
bibstyle=authoryear,
citestyle=authoryear,
maxcitenames=2,
maxbibnames=99
]{biblatex}

\setlength\bibitemsep{1\itemsep} % spacing between entries in references
\addbibresource{references.bib} % .bib file


\usepackage{hyperref} %Must be loaded at the end.
\hypersetup{ %Setup of the reference links.
     colorlinks=false,
     linkcolor=blue,
     citecolor=blue,
     filecolor=magenta,
     urlcolor=blue
}

\usepackage[nameinlink,capitalize]{cleveref} %For the command \cref, load after hyperref.

%----------------------------------------------------------------------------------------
% GENERAL  - SETUP
%----------------------------------------------------------------------------------------
%Define some reasonable margins.
\setlength{\textwidth}{6.6in}
\setlength{\textheight}{8.8in}
\setlength{\topmargin}{-0.1in}
\setlength{\oddsidemargin}{0in}
\setlength{\parskip}{1mm}

\setlength{\parindent}{0cm} %Uncomment this if you don't want to have indents.

%Read just the numbering.
\counterwithout{footnote}{chapter}
\numberwithin{equation}{chapter}

\allowdisplaybreaks[1] %Page breaks of equations are allowed, but avoided if possible. 2-4 more relaxed.

%----------------------------------------------------------------------------------------
% CUSTOM COMMANDS/ENVIRONMENTS
%----------------------------------------------------------------------------------------
%New command for the differential d to have an ordinary d.
\makeatletter
  \newcommand{\ud}{\mathrm{d}}
\makeatother

%Declare Definitions, Theorems etc.(ENVIRONMENTS)
\declaretheorem[style=definition,qed=\(\blacktriangleleft\), numberwithin=chapter]{remark} %additional options; numberwithin=,..., see 'Thmtools' Users’ Guide
\declaretheorem[style=definition,qed=\(\triangle\),numberwithin=chapter]{definition}
\newtheorem{ass}{Assumption}[chapter]
\newtheorem{prop}{Proposition}[chapter]
\newtheorem{lemma}{Lemma}[chapter]
\declaretheorem[style=definition,qed=\(\perp\),numberwithin=chapter]{example}
\newtheorem{theorem}{Theorem}[chapter]
\newtheorem{coroll}{Corollary}[chapter]

%----------------------------------------------------------------------------------------
% TITLE PAGE -  Creates titlepage command
%----------------------------------------------------------------------------------------
%New command for the UZH logo. Used within big \titleGP
\newcommand*{\uzhlogo}{\includegraphics{Graphics/uzh_logo_e_pos.pdf}}


\newcommand*{\titleGP}{\begingroup %Create the command for including the title page in the document.

\centering %Center all text.

% Include University logo
\vspace*{\baselineskip} %White space at the top of the page.
\ifQF
\uzhlogo\hspace{120pt}\ethlogo\\[2\baselineskip] %University Logo.
\else
\uzhlogo\\[2\baselineskip] %University Logo.


% Create title 'box'
\rule{\textwidth}{1.6pt}\vspace*{-\baselineskip}\vspace*{2pt} %Thick horizontal line.
\rule{\textwidth}{0.4pt}\\[\baselineskip] %Thin horizontal line.

% TITLE
{\LARGE Momentum Investing: What has been the best lookback period for a momentum strategy on the set of us industry portfolios?}\\[0.2\baselineskip] %Title.

\rule{\textwidth}{0.4pt}\vspace*{-\baselineskip}\vspace{3.2pt} %Thin horizontal line.
\rule{\textwidth}{1.6pt}\\[2\baselineskip] %Thick horizontal line.
\scshape %Small caps.

% Thesis
\thesis\\
 \par
\fi

\vspace*{2\baselineskip}

% AUTHOR BLOCK - ADAPT
Authors\\
{\Large Aleksandar Kuljanin  \\ Premton Sherifi  \\ Naser Avduli   \\ Luca Müller \\ [5pt]
 }


\vfill

{\scshape Date of Submission: December 11, 2024} \\[0.3\baselineskip]

\endgroup} % This is the end of the command


%----------------------------------------------------------------------------------------
% HEADER/FOOTER
%----------------------------------------------------------------------------------------
% Special header and footer style for the Executive summary and Task Assignment section.
\fancypagestyle{firststyle}{%
  \fancyhf{}%
  \renewcommand{\headrulewidth}{0pt}
  \fancyfoot[C]{\thepage}
}

%Customize headers and footers. - ADAPT
\pagestyle{fancy}
\fancyhead[R]{\thepage}
\fancyhead[L]{\rightmark}
\fancyfoot[L]{} % NAME
\fancyfoot[C]{}
\fancyfoot[R]{Momentum Strategy Analysis} % RUNNING TITLE
\setlength{\headheight}{13.6pt}

%----------------------------------------------------------------------------------------
% SIGNATURE setup
%----------------------------------------------------------------------------------------
%Define the signature line with dots. (create \dotbox command)
\NewDocumentCommand\dotbox{o O{.5\linewidth} m O{3ex} O{\linewidth}}
{
  \begin{minipage}{7cm}
    \makebox[7cm][l]{\,.\dotfill}
    \\
    \makebox[7cm][l]{\,#3}
  \end{minipage}
}